\begin{titlepage}
\phantomsection
\label{sec:oc133-title-page}
\thispagestyle{empty}
\definecolor{ocTitleBlue}{HTML}{173B57}
\definecolor{ocTitleGold}{HTML}{A77D2A}
\begin{tikzpicture}[remember picture,overlay]
\draw[ocTitleBlue,line width=0.75pt]
  ([xshift=1.35cm,yshift=-1.35cm]current page.north west)
  rectangle
  ([xshift=-1.35cm,yshift=1.35cm]current page.south east);
\draw[ocTitleGold,line width=0.35pt]
  ([xshift=1.55cm,yshift=-1.55cm]current page.north west)
  rectangle
  ([xshift=-1.55cm,yshift=1.55cm]current page.south east);
\end{tikzpicture}
\begin{center}
\vspace*{1.0cm}
{\sffamily\Large\color{ocTitleBlue}\textsc{Ontology of Continua}}\\[0.38cm]
{\sffamily\Huge\bfseries Core~1.3.3}\\[0.28cm]
{\sffamily\Large Master Monograph}\\[0.75cm]
{\color{ocTitleGold}\rule{0.42\textwidth}{0.5pt}}\\[0.95cm]

{\Large Alexander Yashin}\\[0.2cm]
{\normalsize Independent Researcher, Leipzig/Halle, Germany}\\[0.2cm]
{\normalsize \href{https://orcid.org/0009-0008-6166-0914}{ORCID 0009-0008-6166-0914}}\\[1.0cm]

{\normalsize Version v1.3.3}\\[0.3cm]
{\normalsize Manuscript revision date: 4 May 2026}\\[0.3cm]
{\normalsize Concept DOI: \href{https://doi.org/10.5281/zenodo.17899134}{10.5281/zenodo.17899134}}\\[0.9cm]

\vfill

% DEDICATION_REQUIRED_DO_NOT_REMOVE
% OC_CORE_PUBLICATION_DEDICATION
{\small\itshape Dedicated to my dear wife Maria, without whom this work would have been impossible.}\\[0.8cm]
\end{center}
\end{titlepage}

\clearpage
\noindent\textbf{Acknowledgements.}
The author thanks the reviewers and critics whose questions, objections, and
suggestions contributed to the development of the Ontology of Continua and to
the refinement of this release. Their criticism helped sharpen the claim
boundaries, improve the reader route, expose weak presentation choices, and
force clearer separation between model claims, proof routes, numerical evidence,
prior-art comparison, and publication form. Acknowledgement records review
pressure and intellectual contribution to the development of the work; it does
not imply authorship, endorsement, publication approval, responsibility for the
theory, or agreement with any claim promoted in the manuscript.
\begin{itemize}
    \item G.~V.~Apostolov.
    \item Eduard Fadeev.
    \item Gennady Alekseevich Nosov.
    \item Sergey Shpadyrev.
    \item Stanislav Tsukrov.
\end{itemize}

\clearpage
\begin{abstract}
Modern systems rarely respect the borders by which modern institutions divide
knowledge. A climate model, an AI platform, a biological organism, a market, an
enterprise architecture, and a scientific theory all involve persistence,
constraint, boundary, transition, failure, repair, and cross-level dependence,
yet they are usually described in mutually incompatible vocabularies. The
Ontology of Continua (OC) begins from that practical and scientific difficulty:
it asks whether heterogeneous systems can be compared through a shared
structural language without erasing the local knowledge that makes each domain
serious.

OC Core~1.3.3 is the bounded external-review release of that model core. It
presents a typed account of continuants, realizations, liveness, death,
residue, rebirth, morphisms, generalized boundaries, hybrid operators, cycle
modes, historical and effective dimension, and K-level witnesses under explicit
assumptions. Its scientific task is not to replace physics, chemistry, biology,
cognition, systems engineering, economics, or philosophy with a slogan. Its
task is to give readers a precise object whose definitions, proof routes,
finite semantic checks, figures, tables, numerical anchors, comparator
boundaries, and failure conditions can be inspected.

The manuscript is written as one continuous scientific monograph rather than as
a prior volume followed by an update packet. It first gives the reader a reason
to care about a continuum ontology, then introduces the elementary picture of a
continuum, the nesting of continua inside continua, and the K-level hierarchy,
and only then moves into the formal core, evidence surfaces, proof routes,
domain projections, limitations, and appendices. The ambition is didactic as
well as formal: the reader should be able to see what the theory is for before
being asked to evaluate its notation.

The current maturity of OC Core is deliberately stated as bounded. The release
is mature enough to be read as a coherent scientific manuscript and to be
checked against formal and executable artifacts. It is not a declaration of
complete scientific coverage, complete numerical closure, or unbounded
comparative victory over contemporary science. Where a theorem route, proof
sheet, Lean declaration, finite semantic case, replay row, comparator row, or
negative-control route supports a claim, the claim is promoted. Where that
support is incomplete, the text records a limitation or future obligation
instead of widening the public wording.

The manuscript contains the full monograph body, a compact article route, a
methods and reproducibility route, reviewer-response material, proof and
formalization sections, finite-model and numerical evidence anchors, figure and
table material, source-trace appendices, and practical-use boundaries. Figures
and tables are retained as scientific material; formulas and theorem statements
are retained in the mathematical spine; machine-readable registers remain in
the evidence package where they can be audited without becoming the narrative
structure of the manuscript.

The practical value of the release is a disciplined compression without loss of
meaning. If OC is useful, it should help a scientific reviewer, a systems
theorist, an AI builder, an enterprise architect, or an evidence auditor
translate a complex system into a common ontology while preserving the
differences that matter. The principal limitation is part of the method:
future releases must add evidence, mechanization, comparator work, or domain
validation before they widen the public claim surface. Version~1.3.3 therefore
offers a publication-grade basis for external review, not a claim that every
future scientific obligation has already been discharged.
\end{abstract}

\clearpage
\section*{Version 1.3.3 Release Delta}
\addcontentsline{toc}{section}{Version 1.3.3 Release Delta}
The OC Core release sequence is part of a continuing scientific program rather
than a series of isolated archive events. As the model is tested against formal
criticism, domain examples, reproducibility checks, and external review, the
public manuscript is expected to become more precise, more explicit about its
limits, and more useful to readers outside the author's immediate working
context. A new release is warranted only when there is a meaningful scientific
or editorial delta: a clarified definition, stronger proof route, better
evidence boundary, improved reader pathway, repaired notation, new domain
projection, or more honest comparator position.

This policy is a commitment to system without pretending that research can be
made mechanical. The author intends to make future public updates regular
enough for readers, reviewers, and auditors to follow the development of the
model, while preserving the difference between a public scientific result and
private working material. No internal route, unpublished process detail, or
control-plane record is promoted as reader-facing evidence merely because it
helped produce the release.

Version~1.3.3 is the release in which the OC Core corpus is reorganized from
scattered source witnesses and process records into a reviewable scientific
package. The visible delta is not merely a new archive number: the release
strengthens the typed foundation, makes the claim boundary explicit, integrates
the theorem route with public proof sheets, and separates reader-facing prose
from machine evidence.

The model delta includes clearer treatment of carriers, realizations, liveness,
death, residue, rebirth, morphisms, generalized boundaries, hybrid operators,
cycle modes, historical and effective dimension, and K-level witnesses. These
additions are presented as a bounded model-core grammar rather than as an
unrestricted theory of every phenomenon.

The verification delta adds and consolidates a Lean-checked subset, structured
proof sheets, finite semantic checks, finite witness interpretation, bounded
target-blind replay QA, comparator and prior-art positioning, negative-control
and falsifier language, and adversarial-review response material. Where a
stronger claim is not yet earned, the release records the limitation instead of
hiding it inside a source-control record.

The editorial delta is equally important. Figures, tables, formulas, numeric
anchors, appendices, and evidence routes are kept in the publication corpus;
machine coverage maps remain coverage and trace data only. The public
documents are therefore built from the human manuscript hierarchy and curated
payload sources, while source bindings and machine evidence indexes stay in the
review manifest.

For a returning reader, the practical point is this: 1.3.3 should be read as a
clarification and consolidation release. It does not claim final closure. It
improves the public surface through which the model can be criticized, taught,
checked, and extended.

\clearpage
\section*{Reader Routes}
\addcontentsline{toc}{section}{Reader Routes}
Dear readers, systems theory is of interest to many communities, but rarely in
exactly the same way. A formal reviewer, a philosopher of systems, an applied
architect, an AI engineer, a security specialist, and an executive reader may
all ask legitimate questions of the same manuscript. OC Core~1.3.3 is therefore
designed as a deliberately generous scientific reading surface: it aims to
disclose the Ontology of Continua as an applied model of system architecture
while giving different readers a courteous route into the material.

The monograph is intentionally somewhat redundant. It first motivates the
problem, then introduces the basic picture of a continuum, then gives the
formal grammar, and only after that asks the reader to inspect proofs, numbers,
domain routes, and appendices. This structure is meant to serve exacting
readers from different scientific and economic contexts without forcing every
reader to begin at the same level of abstraction.

\paragraph{Scientific reviewers and formal critics.}
You may wish to start with the Abstract, Release Delta, Part~I, the formal core
in Part~III, the theorem/proof route, the Lean subset, finite semantic
witnesses, proof machinery appendix, and falsifiability sections. This route is
meant to make the work attackable in the best sense: every promoted claim
should expose assumptions, formulas, proof or executable evidence, comparator
boundary, and reopening condition.

\paragraph{Systems theorists, philosophers, and interested theoretical readers.}
You may prefer to start with the continuum problem, historical and systems-
theory motivation, the K-level primer, lifecycle and boundary chapters,
prior-art comparison, and limitations before opening numeric evidence. This
path asks what OC inherits, what it reorganizes, where its residual delta
begins, and where predecessors or parallel branches already carry part of the
claim.

\paragraph{Applied architects of complex systems.}
Engineers, enterprise architects, AI builders, security specialists, and
applied researchers may find it useful to start from practical examples,
domain projections, figures, tables, and worked routes before returning to the
formal core. The model chapters can then be read as a design grammar for
carriers, realizations, boundaries, operators, update/flow separation, evidence
classes, and claim boundaries in AI systems, evaluation pipelines, simulations,
knowledge models, and formalized domain architectures.

\paragraph{CIO, CEO, enterprise architecture, and strategy readers.}
Readers approaching the manuscript from institutional or strategic work may
sensibly read the release delta, the practical utility chapter, the
model-comparison appendix, and the final conclusion before investing in proof
details. This route asks what the model is for, what it can support today, what
remains too immature for operational commitment, and how evidence classes
affect research or enterprise decisions.

\paragraph{Reproducibility auditors, evidence reviewers, and benchmark readers.}
Readers responsible for verification may go directly to the methods and
evidence chapters, target-blind replay QA, numeric tables, machine-readable
table reference, audit trail, source-trace appendix, and public evidence
package manifest. This route asks whether formula, input snapshot, comparator,
residual or uncertainty, negative control, falsifier, replay hash, and failure
interpretation are present.

\paragraph{Structure of the monograph.}
Part~I explains why a continuum ontology is being proposed at all. Part~II
introduces continua, nesting, and K-levels in reader-facing language. Part~III
states the formal core. Part~IV gathers evidence, proof, falsifiability,
predictions, and reproducibility routes. Part~V presents domain and practical
routes. Part~VI closes with limitations, prior art, and citation/repository
guidance. The appendices carry proof, evidence, table, and source-trace support.

\paragraph{Reading rule shared by all groups.}
The author asks all readers to treat limitations as part of the claim, not as a
footnote. A statement is promoted only where its assumptions, proof or evidence
class, comparator boundary, formulas, appendices, numeric or formal anchor,
figure/table explanation, and reopening condition are visible. Claims about
complete scientific coverage, complete numerical closure, or unbounded
comparative victory over contemporary science are not promoted by this release.

\clearpage
